% ==============================================================================
% SECTION 8: PRBS DATA GENERATION
% ==============================================================================
\section{Data Generation}

\begin{frame}{Data generation}
    \framesubtitle{PRBS excitation in the 3D drone workspace}

    \begin{columns}[T,onlytextwidth]

        \begin{column}{0.48\textwidth}
            \begin{block}{Objective}
                Generate sufficiently rich input/output trajectories to
                identify a nonlinear dynamic model of the drone.
                \begin{itemize}
                    \item Excite the four actuator commands around hover
                    \item Capture the coupled 3D dynamics (12 variables)
                    \item Store synchronized inputs and measured outputs
                \end{itemize}
            \end{block}

            \begin{block}{Simulation setup}
                \begin{itemize}
                    \item Sampling time: $T_s = 0.05\,\mathrm{s}$
                    \item $N = 2000$ samples per trajectory
                    \item $3000$ independent experiments
                    \item Initial hover command: $U_{\mathrm{hover}} = 0.55$
                \end{itemize}
            \end{block}
        \end{column}

        \begin{column}{0.48\textwidth}
            \begin{block}{APRBS modes}
                We don't directly generate u inputs (unusable data).

                Four independently generated piecewise-constant modes:
                \[
                    \delta = [\delta_T,\delta_r,\delta_p,\delta_y]
                \]
                
                Each level is sampled uniformly and held for a random duration
                between $1$ and $3$ seconds.
            \end{block}

            \begin{exampleblock}{Initial phase}
                The excitation starts after $t_{\mathrm{start}}=2\,\mathrm{s}$,
                allowing the drone to stabilize at hover.
            \end{exampleblock}
        \end{column}
    \end{columns}
\end{frame}

\begin{frame}{Data generation}
    \framesubtitle{From modal excitation to motor commands}
    \begin{columns}[T,onlytextwidth]
        \begin{column}{0.48\textwidth}
            \begin{block}{Mixer matrix}
                The four modal perturbations are mixed into the actuator
                commands:
                \[
                    U = U_{\mathrm{hover}} +
                    \delta
                    \begin{bmatrix}
                    1&1&-1&1\\
                    1&-1&-1&-1\\
                    1&-1&1&1\\
                    1&1&1&-1
                    \end{bmatrix}^{\!T}
                \]
                The resulting commands are clipped to the allowed range.
            \end{block}
        \end{column}

        \begin{column}{0.48\textwidth}
            \begin{block}{Input/output data}
                \begin{itemize}
                    \item Inputs: $U_1,\ldots,U_4$
                    \item Outputs: angles, angular velocities,
                          accelerations and angular accelerations
                    \item One trajectory is stored as a sequence of $N$
                          time steps
                    \item Different random seeds improve excitation diversity
                \end{itemize}
            \end{block}

        \end{column}
    \end{columns}
\end{frame}
